\documentclass[11pt,a4paper]{article}
\usepackage[utf8]{inputenc}
\usepackage[T1]{fontenc}
\usepackage[brazil]{babel}
\usepackage{lmodern}
\usepackage[margin=2.3cm]{geometry}
\usepackage{xcolor}
\usepackage{listings}
\usepackage{booktabs}
\usepackage{tabularx}
\usepackage{parskip}
\usepackage[hidelinks]{hyperref}

\definecolor{fundo}{gray}{0.96}
\definecolor{coment}{rgb}{0.35,0.45,0.35}
\definecolor{chave}{rgb}{0.55,0.10,0.35}
\definecolor{texto}{rgb}{0.15,0.35,0.60}

\lstdefinelanguage{Rust}{
  morekeywords={fn,let,mut,match,if,else,for,in,return,struct,enum,impl,use,
    where,pub,mod,self,Self,Some,None,Ok,Err,true,false,continue,as,break},
  morekeywords=[2]{Result,Option,String,str,f64,usize,File,BufReader,io},
  sensitive=true,
  morecomment=[l]{//},
  morestring=[b]",
}
\lstset{
  language=Rust,
  basicstyle=\ttfamily\small,
  keywordstyle=\color{chave}\bfseries,
  keywordstyle=[2]\color{texto},
  commentstyle=\color{coment}\itshape,
  stringstyle=\color{texto},
  backgroundcolor=\color{fundo},
  frame=single, rulecolor=\color{gray!40},
  breaklines=true, columns=fullflexible, keepspaces=true,
  showstringspaces=false, numbers=left, numberstyle=\tiny\color{gray},
  xleftmargin=1.5em, framexleftmargin=1.5em,
  literate={á}{{\'a}}1 {é}{{\'e}}1 {í}{{\'i}}1 {ó}{{\'o}}1 {ú}{{\'u}}1
           {ã}{{\~a}}1 {õ}{{\~o}}1 {ç}{{\c c}}1 {ê}{{\^e}}1 {â}{{\^a}}1
           {ô}{{\^o}}1 {à}{{\`a}}1 {É}{{\'E}}1 {Ç}{{\c C}}1 {—}{{---}}1
}
\lstdefinestyle{saida}{language={},numbers=none,basicstyle=\ttfamily\footnotesize,
  backgroundcolor=\color{white},xleftmargin=0pt,framexleftmargin=0pt}

\newcommand{\code}[1]{\texttt{#1}}

\title{Estudo Dirigido --- Tratamento de Exceções em Rust}
\author{Conceitos de Linguagens de Programação\\[2pt]
  \small Aluno: Gustavo Marcos Brendler}
\date{Outubro de 2026}

\begin{document}
\sloppy
\maketitle

\section{Implementação}

O programa recebe o nome do arquivo pela linha de comando
(\code{cargo run -- dados/correto.txt}), lê cada linha no formato
\code{nome;nota1;nota2}, calcula a média e imprime o resultado.
Os erros foram divididos em dois níveis:

\begin{itemize}
  \item \textbf{Erros de E/S} (abrir ou ler o arquivo): representados por
    \code{std::io::Error} e propagados com \code{?} até \code{main}, que decide
    a mensagem e encerra o programa (não há como continuar sem o arquivo).
  \item \textbf{Erros de registro} (campo ausente, nota inválida, campos a mais):
    representados por um tipo próprio, o \code{enum ErroRegistro}, e tratados
    linha a linha com \code{match}, de modo que um registro inválido é apenas
    ignorado e os seguintes continuam sendo processados.
\end{itemize}

\lstinputlisting[firstline=9]{src/main.rs}

\newpage
\section{Execução dos casos de teste}

\lstinputlisting[style=saida]{dados/saidas.txt}

Mensagens de erro são escritas na saída de erro (\code{eprintln!}) e o
resumo final na saída padrão. Além dos quatro casos, o programa inclui
testes unitários (\code{cargo test}), todos aprovados.

\section{Questões}

\subsection*{1. Como cada problema foi representado e tratado?}

\begin{tabularx}{\linewidth}{@{}p{2.6cm}XX@{}}
\toprule
\textbf{Problema} & \textbf{Representação} & \textbf{Tratamento} \\
\midrule
Arquivo inexistente &
\code{File::open} devolve \code{Err(io::Error)} com
\code{kind() == ErrorKind::NotFound}. &
Propagado com \code{?} em \code{processar\_arquivo}; tratado em \code{main}
por um \code{match} com guarda que testa o \code{kind()}; exibe mensagem e
encerra com código 1. \\
\addlinespace
Valor numérico inválido &
\code{"dez".parse::<f64>()} devolve \code{Err(ParseFloatError)}, convertido
por \code{map\_err} em \code{ErroRegistro::NotaInvalida}. &
Propagado com \code{?} de \code{converter\_nota} até \code{analisar\_linha} e
desta até o laço; ali um \code{match} registra o erro e o laço segue para a
próxima linha. \\
\addlinespace
Campo ausente &
\code{split(';').next()} devolve \code{None} (um \code{Option}), que é
convertido em \code{Err(ErroRegistro::CampoAusente)}. &
Mesmo caminho do caso anterior: propagado por \code{?} e tratado no
\code{match} do laço, sem interromper o processamento. \\
\bottomrule
\end{tabularx}

Em Java essas situações corresponderiam, respectivamente, a
\code{FileNotFoundException}, \code{NumberFormatException} e
\code{ArrayIndexOutOfBoundsException} (ou uma exceção própria). Em Rust
nenhuma delas é uma exceção: todas são \emph{valores} comuns de retorno.
O \emph{tipo} da função (\code{Result<Estudante, ErroRegistro>},
\code{io::Result<...>}) anuncia quais erros podem ocorrer, e o compilador não
permite usar o valor obtido sem antes considerar o caso de erro. Vale notar que a ausência de campo
nasce como \code{Option} (``não há valor'', sem motivo associado) e só
vira erro quando o programa decide que essa ausência é um problema e lhe
dá um significado (\code{CampoAusente \{ campo: "nota2" \}}).

\subsection*{2. Uma operação que retorna \code{Result<T, E>}}

A função \code{converter\_nota} tem tipo
\code{fn(\&str, \&'static str) -> Result<f64, ErroRegistro>}:

\begin{lstlisting}[numbers=none]
fn converter_nota(texto: &str, campo: &'static str) -> Result<f64, ErroRegistro> {
    texto.parse::<f64>().map_err(|causa| ErroRegistro::NotaInvalida {
        campo, valor: texto.to_string(), causa,
    })
}
\end{lstlisting}

\begin{itemize}
  \item \code{Ok(f64)}: a conversão foi bem-sucedida; o valor carregado é a
    nota numérica (ex.: \code{"8.5"} $\to$ \code{Ok(8.5)}).
  \item \code{Err(ErroRegistro)}: a conversão falhou; o valor carregado
    descreve a falha --- qual campo, qual texto e a causa original
    (ex.: \code{"dez"} $\to$ \code{Err(NotaInvalida \{ campo: "nota1",
    valor: "dez", ... \})}).
\end{itemize}

\textbf{Relação com o levantamento de exceção.} Na videoaula, levantar uma
exceção (\code{throw}/\code{raise}) significa sinalizar que uma condição
anormal ocorreu, interrompendo o fluxo normal e transferindo o controle
para um tratador, junto com um objeto que descreve a exceção. Em Rust,
\emph{retornar \code{Err(...)}} cumpre o papel de sinalização: a função
termina e entrega ao chamador um objeto descritivo do erro (o análogo do
objeto de exceção, com ``tipo'' e dados). A diferença é que nada é
desviado automaticamente: o \code{Err} é um valor de retorno normal, que
segue pelo caminho usual de retorno de função. O ``levantamento'' é,
portanto, explícito e local, e o tipo \code{Result} torna a possibilidade
de erro parte da assinatura --- algo parecido com as exceções
\emph{verificadas} de Java (\code{throws}), porém imposto pelo sistema de
tipos e não por uma regra à parte.

\subsection*{3. O operador \code{?}}

O \code{?} é usado em \code{analisar\_linha} e em \code{processar\_arquivo}:

\begin{lstlisting}[numbers=none]
let nota1 = converter_nota(proximo_campo(&mut campos, "nota1")?, "nota1")?;
let arquivo = File::open(caminho)?;
\end{lstlisting}

\begin{itemize}
  \item \textbf{Quando a expressão produz \code{Ok(v)}}: o \code{?}
    ``desembrulha'' o valor e a expressão inteira vale \code{v}. A execução
    continua normalmente na linha seguinte (\code{nota1} recebe o
    \code{f64}; \code{arquivo} recebe o \code{File}).
  \item \textbf{Quando produz \code{Err(e)}}: a função atual retorna
    \emph{imediatamente} \code{Err(From::from(e))} ao seu chamador; o
    restante do corpo não é executado. O \code{From::from} permite converter
    automaticamente o tipo de erro, se houver conversão definida.
\end{itemize}

Ou seja, \code{x?} equivale a:
\begin{lstlisting}[numbers=none]
match x { Ok(v) => v, Err(e) => return Err(From::from(e)) }
\end{lstlisting}

\textbf{Comparação com propagação de exceções.} No modelo tradicional,
quando uma unidade não possui tratador para a exceção levantada, ela é
encerrada e a exceção é \emph{propagada} automaticamente para o chamador
(seguindo a cadeia dinâmica de chamadas), até encontrar um tratador
ou chegar ao topo e abortar o programa. O \code{?} produz o mesmo efeito
observável --- a função é encerrada e o erro ``sobe'' para o chamador ---,
porém com diferenças importantes:
\begin{itemize}
  \item a propagação é \textbf{explícita}: cada ponto em que um erro pode
    sair da função está marcado com \code{?}; em Java, qualquer chamada
    pode silenciosamente propagar uma exceção não verificada;
  \item ela ocorre \textbf{um nível por vez}: cada função intermediária
    precisa declarar um tipo de retorno compatível (\code{Result}) e usar
    \code{?} de novo; não há ``salto'' por várias pilhas de uma só vez;
  \item não existe desempilhamento especial: é um \code{return} comum.
\end{itemize}
No programa, a nota inválida ``dez'' é propagada de \code{converter\_nota}
$\to$ \code{analisar\_linha} $\to$ laço em \code{processar\_arquivo}, onde
é tratada; já o \code{NotFound} vai de \code{File::open} $\to$
\code{processar\_arquivo} $\to$ \code{main}.

\subsection*{4. Tratamento explícito com \code{match} versus \code{try-catch}}

\begin{minipage}[t]{0.49\linewidth}
\textbf{Rust}
\begin{lstlisting}[numbers=none,basicstyle=\ttfamily\scriptsize]
match analisar_linha(&linha) {
  Ok(estudante) => {
    println!("{} média: {:.2}", ...);
  }
  Err(erro) => {
    eprintln!("[linha {}] ... {}",
              numero, erro);
    invalidos += 1;
  }
}
\end{lstlisting}
\end{minipage}\hfill
\begin{minipage}[t]{0.49\linewidth}
\textbf{Java (equivalente)}
\begin{lstlisting}[numbers=none,language=Java,basicstyle=\ttfamily\scriptsize]
try {
  Estudante e = analisarLinha(linha);
  System.out.printf(...);
} catch (NumberFormatException ex) {
  System.err.println(...);
  invalidos++;
} catch (CampoAusenteException ex) {
  System.err.println(...);
  invalidos++;
}
\end{lstlisting}
\end{minipage}

\textbf{Onde está o tratador e como é associado ao erro.}
\begin{itemize}
  \item \textbf{Java:} o tratador é a cláusula \code{catch} associada a um
    bloco \code{try}. A associação é \emph{dinâmica} e feita pelo
    \emph{tipo} da exceção: quando algo é lançado em qualquer ponto do
    bloco protegido (inclusive dentro de métodos chamados por ele), a
    execução é desviada e a JVM procura, entre os \code{catch} desse
    \code{try} e depois dos \code{try} envolventes na pilha de chamadas,
    o primeiro cujo tipo é compatível. O código protegido e o código de
    tratamento ficam separados.
  \item \textbf{Rust:} o ``tratador'' é o braço \code{Err(erro) => ...} do
    \code{match}, e ele se associa a \emph{um valor específico} --- o
    resultado daquela única chamada \code{analisar\_linha(\&linha)}. A
    associação é \emph{estática e léxica}: é o próprio programador que
    inspeciona o valor retornado, no ponto exato da chamada. A seleção entre
    tipos de erro se faz por casamento de padrões (ex.:
    \code{Err(e) if e.kind() == NotFound} em \code{main}, ou
    \code{Err(ErroRegistro::CampoAusente \{..\})}), e o compilador exige
    que o \code{match} seja exaustivo, ou seja, que \code{Ok} e \code{Err}
    sejam ambos considerados.
\end{itemize}
Em ambos os casos, após o tratamento a execução continua na instrução
seguinte ao \code{try}/\code{match} --- o modelo de \emph{terminação}, não
de retomada: o registro com erro não é reprocessado, o laço apenas segue
para a próxima linha.

\subsection*{5. Rust possui tratamento de exceções no sentido da disciplina?}

\textbf{Não, para erros recuperáveis.} No sentido apresentado na disciplina
--- um mecanismo da linguagem no qual uma exceção é \emph{levantada},
desvia implicitamente o fluxo de controle e é \emph{propagada
automaticamente} pela cadeia de chamadas até um tratador associado
dinamicamente --- Rust não oferece esse mecanismo para erros esperados.
Esses erros são tratados como \textbf{valores} (\code{Result} e
\code{Option}), com apoio de \code{match}, do operador \code{?} e de
combinadores (\code{map\_err}, \code{ok\_or}, \code{unwrap\_or}, \dots).

O que mais se aproxima de exceção em Rust é o \code{panic!}: ele interrompe
o fluxo e desempilha a pilha (\emph{unwinding}) executando destrutores, e
pode até ser capturado com \code{std::panic::catch\_unwind}. Porém, ele se
destina a erros \emph{irrecuperáveis} (bugs, violação de invariantes),
pode ser configurado para simplesmente abortar (\code{panic = "abort"}) e
não é a forma idiomática de sinalizar erros como os deste estudo.

\begin{tabularx}{\linewidth}{@{}p{3.1cm}XX@{}}
\toprule
\textbf{Aspecto} & \textbf{Modelo tradicional (Java, C++, Python)} & \textbf{Rust} \\
\midrule
Sinalização & \code{throw} de um objeto exceção & \code{return Err(valor)} (ou \code{None}) \\
\addlinespace
Representação do erro & Objeto de uma hierarquia de classes de exceção & Valor de qualquer tipo, geralmente um \code{enum} \\
\addlinespace
Fluxo de controle & Desvio implícito, não local & Retorno comum, local e explícito \\
\addlinespace
Propagação & Automática pela pilha até achar tratador & Explícita com \code{?}, um nível por vez \\
\addlinespace
Associação ao tratador & Dinâmica, pelo tipo, ao \code{catch} mais próximo & Léxica, por \code{match} sobre o valor retornado \\
\addlinespace
Erro não tratado & Termina o programa (ou a thread) & Não passa despercebido: \code{Result} ignorado gera aviso do compilador, e para obter o \code{T} é preciso decidir o que fazer com o \code{Err} (\code{match}, \code{?}, \code{unwrap}\dots) \\
\addlinespace
Visibilidade na assinatura & Só exceções verificadas (Java) & Sempre, no tipo de retorno \\
\addlinespace
Continuação & Terminação & Terminação (a função que falha não é retomada) \\
\addlinespace
Erros irrecuperáveis & Mesma construção (\code{Error}, \code{RuntimeException}) & Mecanismo separado: \code{panic!} \\
\bottomrule
\end{tabularx}

\textbf{Semelhanças:} ambos separam o resultado normal do resultado de
erro; ambos carregam informação descritiva sobre a falha (tipo e dados);
ambos permitem tratar o erro no local ou repassá-lo ao chamador; e ambos
seguem o modelo de terminação, em que a unidade que detectou o erro não é
retomada.

\textbf{Diferenças:} em Rust, o erro é um valor comum e não um evento de
controle; a propagação é explícita (\code{?}) e não automática; a
associação com o tratador é léxica e verificada pelo compilador
(exaustividade do \code{match}), o que impede que um erro seja esquecido;
e não há custo nem complexidade de desempilhamento para o caminho de erro.
Em troca, o código fica um pouco mais verboso e cada função intermediária
precisa declarar o tipo de erro que pode devolver --- custo que o
operador \code{?} e a conversão via \code{From} reduzem bastante.

Em resumo, Rust \emph{resolve o mesmo problema} que o tratamento de
exceções resolve (detectar, sinalizar, propagar e tratar situações
anormais), mas por meio do sistema de tipos, e não por um mecanismo de
desvio de controle. Seu único mecanismo semelhante a exceções, o
\code{panic!}, é reservado para falhas irrecuperáveis.

\end{document}
